% !PRINTED 18
\weberpage{18}

If we therefore set $\mu = \frac{1}{12}$, the equations (19)
pass over into the identities $g_2=g_2$, $g_3=g_3$.

If we then set $y=1$, the transformation (18) gives

\weq{20}{\dfrac{2\,dx}{\sqrt{4\,x^3-g_2\,x-g_3}} = \dfrac{dz}{\sqrt{4\,z^3-g_2\,z-g_3}}\,.}

% [sic] Printed "according to (16)"; with \mu=1/12, the displayed substitution follows from (17).
This transformation is mediated, according to (16), by

\[ z = \dfrac{-H}{48\,f} \]

But from (14) it follows that:

\[ H = -48\left[(x^2+\tfrac{1}{4}g_2)^2 + 2\,g_3\,x\right], \]

and consequently we obtain

\weq{21}{z = \dfrac{(x^2+\tfrac{1}{4}g_2)^2+2\,g_3\,x}{4\,x^3-g_2\,x-g_3}\,.}

There follows further from (14)

\[ T = 64\,x^6 - 80\,g_2\,x^4 - 320\,g_3\,x^3 - 20\,g_2^2\,x^2 - 16\,g_2\,g_3\,x + g_2^3 - 32\,g_3^2 \]

and then, according to (16),

\weq{22}{4\sqrt{4\,x^3-g_2\,x-g_3}^{\,3}\ \sqrt{4\,z^3-g_2\,z-g_3} = T\,.}

\gesperrt{By (20), (21), (22) the multiplication of the elliptic
 differential by 2 is accomplished.} Not only is
$z$ thereby represented as a single-valued (rational) function of $x$,
but one square root is also represented uniquely by the other;
i.\,e., a \gesperrt{point} $z$ is uniquely assigned to each \gesperrt{point} $x$.

\begin{center}
§\,6. \gesperrt{Elliptic Curves.}
\end{center}

Every algebraic dependence between two variables
$x$, $y$ is expressed by an equation of the form

\weq{1}{F(x,y) = 0\,,}

in which $F(x,y)$ denotes an integral function of the two variables $x$, $y$.
If one regards $x$ and $y$ as Cartesian coordinates
of a point in the plane, then (1) is the equation of a curve
of degree $n$, provided $F$, with respect to $x$ and $y$ taken together,
is of the $n$th dimension. This curve is then the geometric
 image of the algebraic dependence, where, however,
imaginary points must also be taken into account.
